\documentclass[border = 1mm]{standalone}
% Caption:
% — Fourier Transform.
% Mathematical definition using mathematics convention.
%
% Autor:
% R. Sánchez-Martínez, Jun 2026.
%
% Figure published at:
% instagram.com/ars_mathematica
% Code available at:
% github.com/ro-smtz/ars-mathematica-content/tree/main/figures
%
% Generated with TikZ and compiled preferable with XeLaTeX.

%
\usepackage[utf8]{inputenc}
\usepackage[T1]{fontenc}
\usepackage{microtype}
%
\usepackage{amsmath,amssymb,amsfonts}
\usepackage{physics}
%
\usepackage{xcolor}
\usepackage{tikz}
\usetikzlibrary{
	babel,
	math,calc
}
\usepackage{tikz-feynman}

\usepackage{fontspec}
\usepackage[default]{fontsetup}

\begin{document}
% -- Figure
% Key concepts:
%	— \draw[style] (a) to (b);
%	  draws a line from (a) to (b) with a specific
%     graphic [style].
%   — \draw (a) to node {text} (b);
%	  adds a node containing {text} in the midpoint
%     of (a) and (b).
%   — ... node[anchor = base] {text} ...
%     indicates that the coordinate specification will
%     refer to the node's baseline midpoint, useful
%     for aligning nodes of different heights.
\begin{tikzpicture}[
	line join = round,
	line cap  = round,
	line width = 0.62pt,
]

	% Main box 
	\node[
		inner sep = 0pt,
		minimum width  = 5cm,
		minimum height = 5cm,
		name = N,
	] at (0,0) {};
	
	% Clipping mask, it deletes everything outside the box (N),
	% useful for obtaining well-defined-size drawings
	\clip (N.north west) rectangle (N.south east);

	% ---
	\begin{scope}[
	% use shift to center the picture or placing it
	% wherever you want
		shift = {(0,0)},
	]
		
		% main equation
		\node at (0,0.6) {
		$\displaystyle
		\begin{aligned}
			\mathcal{F}\big[f(t)\big]
			&= \int\limits_{-\infty}^{\infty}\!
			f(t)\, e^{2\pi\mathrm{i}\,\omega t}
			\dd{t} \\
			&= F(\omega)
		\end{aligned}
		$
		};
		
		%
		\node at (0,-1.2) {
		$\displaystyle
			f\;
			\overset{
				\text{\small$\mathcal{F}$}
			}{
				\textcolor{gray!75}{\Longleftrightarrow}
			}
			\; F
		$
		};
	
	\end{scope}
	% ---
	
\end{tikzpicture}
\end{document}